\section{Timing analysis}
\label{sec:analysis}

\begin{inbrief}
This section gives the analyses that the method calls, with the arguments
that make them sound: response times on a core (\cref{sec:rta}), an
optimal priority assignment that respects precedence (\cref{sec:audsley}),
Network Calculus bounds for CBS, TAS, ATS and FRER (\cref{sec:nc}), the
end-to-end bounds of a task-chain (\cref{sec:e2e}), the critical scaling
factor (\cref{sec:sensitivity}) and a lower bound on the number of cores
(\cref{sec:lower-bound}).
\end{inbrief}

\subsection{Response times on a core}
\label{sec:rta}

Let $p$ be a core, $\hp(\tau)$ the tasks of $p$ with a higher priority
than $\tau$ (the LET task included), and $C_j=C_j(p)$ the WCETs of
\cref{eq:wcet}, contention included. The WCRT of $\tau$ is the least fixed
point of \cite{joseph1986finding,audsley1993applying}
\begin{equation}
R=C_\tau+\sum_{j\in\hp(\tau)}\ceil{\frac{R}{T_j}}C_j .
\label{eq:rta}
\end{equation}
The right-hand side $F(R)$ is non-decreasing in $R$. If $R^{(0)}\le R^\ast$
and $R^{(0)}\le F(R^{(0)})$, the iterates $R^{(i+1)}=F(R^{(i)})$ are
non-decreasing integers bounded by the least fixed point $R^\ast$, so they
reach it in finitely many steps; both starting values of \cref{lem:start}
satisfy the two conditions. The iteration stops early as soon as
$R^{(i)}>T_\tau$ (deadline miss). With contention, \cref{eq:wcet} adds to
each WCET the longest wait of its memory requests, in the spirit of
multicore response-time analysis \cite{davis2018extensible}, with the
round-robin bound of \cite{paolieri2009hardware}; it is safe for any interleaving of the
requests of the other cores.

\begin{lemma}[Starting point]
\label{lem:start}
With $U_{\hp}=\sum_{j\in\hp(\tau)}C_j/T_j<1$, the least fixed point $R^\ast$
satisfies $R^\ast\ge\max\bigl(C_\tau+\sum_{\hp}C_j,\;C_\tau/(1-U_{\hp})\bigr)$.
\end{lemma}
\begin{proof}
Every higher-priority task is released at the critical instant, so
$\lceil R/T_j\rceil\ge1$. Since $\lceil x\rceil\ge x$, also $R^\ast\ge
C_\tau+U_{\hp}R^\ast$.
\end{proof}
In floating point, the second term is computed as $\lfloor C_\tau/(1-
U_{\hp}+10^{-9})\rfloor-1$, which keeps it a valid lower bound despite
rounding.

\begin{lemma}[Grouping by period]
\label{lem:group}
Partition $\hp(\tau)$ into groups $G_1,\dots,G_q$ of equal periods
$T^{(1)},\dots,T^{(q)}$, and let $W_g=\sum_{j\in G_g}C_j$. Then $F(R)=
C_\tau+\sum_{g=1}^{q}\lceil R/T^{(g)}\rceil W_g$, and the iteration produces
the same integer sequence in $O(q)$ instead of $O(|\hp(\tau)|)$
operations per step.
\end{lemma}
\begin{proof}
For $j\in G_g$, $\lceil R/T_j\rceil=\lceil R/T^{(g)}\rceil$; factor it out.
All quantities are integers, so the sums are exact.
\end{proof}
Automotive applications use few distinct periods; the WATERS benchmark has
nine, besides its angle-synchronous tasks \cite{kramer2015waters}. A core with hundreds of tasks therefore costs
a handful of terms per step. To compute all the response times of a core
in priority order, the groups, $\sum_{\hp}C_j$ and $U_{\hp}$ are
accumulated incrementally.

\subsection{Optimal priority assignment under precedence}
\label{sec:audsley}

Audsley's algorithm \cite{audsley2001priority} fills the priority levels
from the lowest. At each level, it gives the level to any task that is
schedulable there with all unassigned tasks above it. It is optimal for
any test in which the schedulability of a task depends on the \emph{set}
of higher-priority tasks, but neither on their order nor on the order of
the lower-priority tasks, and in which a task cannot become unschedulable
when its priority rises \cite{audsley2001priority}. \Cref{eq:rta} is such
a test. Constraint (C6) restricts the admissible orders, so the candidates
for a level must be restricted too.

\begin{definition}[Candidates]
Let $S$ be the unassigned application tasks of core $p$. For $\tau\in S$,
let $\mathrm{next}(\tau)$ be the next task of the same chain and period on
$p$, if any. The \emph{candidates} for the lowest free level are
$\mathcal{C}(S)=\{\tau\in S:\mathrm{next}(\tau)\notin S\}$. A candidate is
\emph{feasible} if its WCRT with $\hp=S\setminus\{\tau\}\cup\{\LET\}$ is at
most its period.
\end{definition}

\begin{algorithm2e}[t]
\caption{\textsf{Priorities}$(p,\textit{robust})$: optimal priority
assignment that respects precedence}
\label{alg:audsley}
\Input{tasks $S_0$ of core $p$ (LET excluded); flag \textit{robust}}
\Output{a priority level for every task; \emph{schedulable}}
$S\gets S_0$\;
\For{$v\gets|S_0|+1$ \KwTo $2$}{
	$\mathcal{C}\gets\{\tau\in S:\mathrm{next}(\tau)\notin S\}$ in order of
	decreasing $T_\tau$, ties by decreasing index\;
	$\textit{best}\gets\bot$; $\textit{bs}\gets-1$; $W\gets C^{\LET}_p+\sum_{\tau\in S}C_\tau$\;
	\ForEach{$\tau\in\mathcal{C}$}{
		\lIf{$(T_\tau-W)/T_\tau<\textit{bs}$}{\textbf{break}\tcp*[f]{bound}}
		$R\gets$ WCRT of $\tau$ below $S\setminus\{\tau\}$ and LET (\cref{lem:start,lem:group})\;
		\lIf{$R>T_\tau$}{\textbf{continue}}
		\lIf{\textbf{not} \textit{robust}}{$\textit{best}\gets\tau$; \textbf{break}}
		\If{$(T_\tau-R)/T_\tau>\textit{bs}$ \textbf{or} (equal \textbf{and} $\tau$ later in its chain)}{
			$\textit{best}\gets\tau$; $\textit{bs}\gets(T_\tau-R)/T_\tau$\;}
	}
	\If{$\textit{best}=\bot$}{
		give the tasks of $S$ the levels $2..v$ in precedence order\;
		\Return \emph{unschedulable}\;}
	$\mathrm{level}(\textit{best})\gets v$; $S\gets S\setminus\{\textit{best}\}$\;
}
LET task $\gets$ level 1 (highest)\;
\Return \emph{schedulable}\;
\end{algorithm2e}

\begin{theorem}[Optimality]
\label{thm:audsley}
If some priority order that satisfies (C6) makes core $p$ schedulable,
then \cref{alg:audsley}, with or without \textit{robust}, returns a
schedulable order that satisfies (C6).
\end{theorem}
\begin{proof}
\emph{(C6) holds by construction}: a task is chosen only when its
successor in the pattern already has a lower level.
\emph{Optimality}, by induction on the levels. Invariant: before the
algorithm fills level $v$, there is a schedulable order that satisfies
(C6) and agrees with the levels already assigned (below $v$). It holds
initially by assumption. Let $O$ be such an order, $S$ the unassigned
tasks (the levels $2..v$ of $O$), and suppose the algorithm picks
$\tau^\ast$ for level $v$. Build $O'$ from $O$ by moving $\tau^\ast$ to
level $v$; the tasks of $S$ that were below it in $O$ move up one level.
(i) $\tau^\ast$ is schedulable at level $v$: its higher-priority set is
$S\setminus\{\tau^\ast\}$ plus the LET task, exactly the set under which the
algorithm checked it, and \cref{eq:rta} depends only on that set. (ii) A
task that moved up loses $\tau^\ast$ from its higher-priority set, and every
other task keeps its set, so no response time grows. (iii) (C6) still holds:
$\tau^\ast$ is a candidate, so no later task of its pattern on $p$ is in
$S$, and the other tasks keep their relative order. So $O'$ satisfies the
invariant for level $v-1$. The invariant also shows that some candidate
is feasible at every level: the task at level $v$ in $O$ is a candidate,
by (C6), and it is schedulable there. Hence, if no candidate is feasible at
some level, no admissible schedulable order exists.
\end{proof}

\paragraph{Robust choice.} Any feasible candidate preserves optimality, but
not the same margin. With \textit{robust}, the level goes to the feasible
candidate with the largest relative slack $(T-R)/T$, after Davis and Burns
\cite{davis2007robust}. The tightest tasks are thereby pushed to higher
levels, which tends to increase the critical scaling factor
(\cref{sec:sensitivity}); this is a heuristic choice, not a guarantee
that the factor is maximal. Because $R\ge W=\sum_{S}C$ (\cref{lem:start}),
$(T-W)/T$ bounds the slack of every candidate. Candidates are visited by
decreasing period, so the bound decreases along the scan, and the scan
stops as soon as it cannot beat the best slack found. A cheaper first-fit
pass (\textit{robust} off) gives the same verdict, so the robust pass runs
only on schedulable cores.

\paragraph{Cost.} The successor links come from one sort of the tasks by
(chain, period, position), in $O(n\log n)$. The eligible tasks are kept as
a bit set in candidate order: choosing a task clears its bit and sets the
bit of its predecessor in the pattern. The period groups and the sums $W$
and $U$ are updated when a task gets a level. A level then costs only the
evaluation of its candidates, $O(c\,\iota\,q)$ for $c$ evaluated candidates,
$\iota$ iterations and $q$ periods.

\subsection{Network Calculus bounds}
\label{sec:nc}

\paragraph{Arrival curves.} Flow $h$ sends one message of $b_h$ bits per
period $T_h$, copied out at a LET instant of the source core, and leaves
when the copy ends. The source core has LET instants at the release of
the producer and every $T^{\LET}_{\mathrm{src}}$ after it. With implicit
communication, the message of a job leaves no earlier than
$T^{\LET}_{\mathrm{src}}$ and no later than
$\lceil R_{\mathrm{prod}}/T^{\LET}_{\mathrm{src}}\rceil
T^{\LET}_{\mathrm{src}}+C^{\LET}_{\mathrm{src}}$ after that release, so its
release jitter is at most $R_{\mathrm{prod}}+C^{\LET}_{\mathrm{src}}\le
J_h=R_{\mathrm{prod}}+T^{\LET}_{\mathrm{src}}$. With strict LET, it leaves
between $T_h$ and $T_h+C^{\LET}_{\mathrm{src}}$ after that release, so
$J_h=C^{\LET}_{\mathrm{src}}$. At its first port the flow is constrained by
the token bucket $\alpha_h(t)=\sigma_h+\rho_h t$, with $\rho_h=b_h/T_h$
and $\sigma_h=b_h+\rho_hJ_h$ \cite{leboudec2001network}. Both use
$R_{\mathrm{prod}}\le T_h$: the WCRT of \cref{eq:rta} is only valid up to
the period, and a producer that misses it would need a busy-period
analysis \cite{lehoczky1990fixed}. Such a flow is given no arrival curve
instead: every port it crosses, and every port these feed, is declared
unbounded, so no bound depends on it.

\paragraph{Port bound.} Let $F(e,x)$ be the flows of class $x$ through
port $e$ and $\sigma_h(e)$ the burst of flow $h$ when it reaches $e$. If
the port gives class $x$ a rate-latency service with rate $R_{e,x}$ and
latency $L_{e,x}$, if $\sum_{F(e,x)}\rho_h<R_{e,x}$, and if the LET load
of the link is within its design limit, then the FIFO aggregate of the
class has the delay bound \cite{leboudec2001network}
\begin{equation}
d_{e,x}=L_{e,x}+\frac{1}{R_{e,x}}\sum_{h\in F(e,x)}\sigma_h(e).
\label{eq:port-delay}
\end{equation}
Otherwise the port is \emph{overloaded} and its flows have no bound. The
three shapers differ by their service curves and by how bursts evolve
along the route.

\paragraph{CBS service.} At a port of rate $r$ with classes $0..X-1$ in
decreasing priority, class $x$ with idle slope $I_x$ receives the
rate-latency service $R=I_x$, $L=L_x$ with
\cite{mohammadpour2018latency,zhao2021latency}
\begin{gather}
L_x=\frac{\cmax_x}{I_x}+\frac{(r-I_x)\,l_x}{r\,I_x},\label{eq:cbs-latency}\\
\cmax_x=\frac{I_x\bigl(r\,\bar l_x+\sum_{j<x}(r-I_j)\,l_j\bigr)}
	{r\,\bigl(r-\sum_{j<x}I_j\bigr)},\label{eq:cbs-credit}
\end{gather}
where $l_j$ is the largest frame of class $j$ at the port and $\bar l_x=
\max(L_{BE},\max_{j>x}l_j)$ the largest frame that can block class $x$
(non-preemptively). With frame preemption, $L_{BE}$ drops to 143 bytes.
The first term of \cref{eq:cbs-latency} is the time to recover the largest
credit. The second accounts for a frame of the class being sent at the
link rate rather than at the idle slope \cite{azua2014complete}.

\paragraph{Idle slopes.} With one class, $I_0$ is a fixed share of $r$
(75\,\%). With several classes, fixed shares waste bandwidth where a class
carries little traffic, so the slopes are set per port. Let $\rho_x=\sum
\rho_h$ over the flows of class $x$ at the port, $\beta_x=\sum b_h$ their
bits, $T^{\min}_x$ their shortest period, and $B=r\sum_x\mathrm{share}_x$
the budget. Then
\begin{equation}
I_x=\rho_x+\Bigl(B-\sum_j\rho_j\Bigr)\frac{\beta_x/T^{\min}_x}
	{\sum_j\beta_j/T^{\min}_j},
\label{eq:slopes}
\end{equation}
if $\sum_j\rho_j\le B$; otherwise the port is overloaded. Each class gets
its own rate, plus a part of the spare budget in proportion to its
burstiness per unit of time, because its delay is about its bursts over
its slope (\cref{eq:port-delay}).

\paragraph{TAS service.} With an exclusive window of length $w_x$ per
gate cycle $T_c$, and guard bands that keep other frames out of the
window, class $x$ is served at the link rate for $\hat w_x=w_x-l_x/r$ per cycle,
because a frame starts only if it ends within the window. A rate-latency
lower bound is $R=r\hat w_x/T_c$, $L=T_c-\hat w_x$ \cite{zhao2018worst}. If $\hat w_x\le0$
the class has no service.

\paragraph{ATS service.} With ATS the classes are served in strict
priority, without preemption. Class $x$ sees the higher classes through
their arrival curves, so it receives the residual rate-latency service
\cite{leboudec2001network}
\begin{equation}
R_{e,x}=r-\sum_{j<x}\rho_j,\qquad
L_{e,x}=\frac{\bar l_x+\sum_{j<x}\sum_{h\in F(e,j)}\sigma_h(e)}{R_{e,x}},
\label{eq:ats}
\end{equation}
where $\bar l_x$ is the largest frame of a lower class or of best effort.

\paragraph{Burst propagation.} A flow that leaves a port with delay bound
$d$ has the output burst $\sigma+\rho d$. With CBS and TAS, along the
route $e_1,\dots,e_n$ of flow $h$,
\begin{equation}
\sigma_h(e_i)=b_h+\rho_h\Bigl(J_h+\sum_{j<i}d_{e_j,x}\Bigr).
\label{eq:burst}
\end{equation}
With ATS, the regulators of every switch hold the frames of each flow until
they conform again to $\alpha_h$. Behind a FIFO system, such regulators do
not increase the worst-case delay of the system plus regulator
(``shaping for free'', \cite{leboudec2018theory}, for both interleaved and
per-flow regulators). Every port of the route therefore sees the source
burst, $\sigma_h(e_i)=b_h+\rho_hJ_h$, and \cref{eq:port-delay} with
\cref{eq:ats} bounds the delay through the port and the regulator that
follows it. The end-to-end bound of a flow is, in all cases,
\begin{equation}
d_h=\sum_{i=1}^{n}d_{e_i,x}+(n-1)\delta .
\label{eq:flow-delay}
\end{equation}

\paragraph{Frame replication.} A replicated flow $h$ has two routes, $P_1$
and $P_2$, with the same first and last ports (\cref{fig:frer}). Both
member streams load the ports of their own route. For route $m$, let
$\bar D_m$ be the delay bound from the source to the elimination: all its
ports but the last, and its $n_m-1$ switches; and let $\underline{d}_m=
(n_m-1)\delta$, a lower bound of the same delay.

\begin{proposition}[Bound with replication]
\label{prop:frer}
After the elimination, flow $h$ is constrained by the token bucket of rate
$\rho_h$ and burst
\begin{equation}
\sigma_h(e_n)=b_h+\rho_h\Bigl(J_h+\max_m\bar D_m-\min_m\underline{d}_m\Bigr),
\label{eq:frer-burst}
\end{equation}
and its delay is at most $d_h=\max_m\bigl(\bar D_m+d_{e_n,x}\bigr)$, also
after the failure of one route.
\end{proposition}
\begin{proof}
Let message $k$ be released at $t_k\in[kT_h,kT_h+J_h]$. Each copy that
reaches the elimination does so within $[t_k+\underline{d}_m,t_k+\bar
D_m]$ for its route $m$, and the elimination forwards the first one. The
forwarded message $k$ therefore leaves the elimination within
$[kT_h+\min_m\underline{d}_m,\,kT_h+J_h+\max_m\bar D_m]$, a window of length
$J_h+\max_m\bar D_m-\min_m\underline{d}_m$. In any interval of length $t$,
at most $1+(t+J_h+\max\bar D_m-\min\underline{d}_m)/T_h$ messages can leave,
of at most $b_h$ bits each, which gives \cref{eq:frer-burst}. This is
corollary~1 of Thomas et al.\ \cite{thomas2022worst} for one
replication-elimination pair. The forwarded copy is never later than the
copy of the slower route, which gives the delay bound. After a failure,
fewer frames cross the ports, and every bound above is monotone in the
traffic, so the bounds still hold.
\end{proof}
With ATS, the regulators of the destination switch act on each member
stream \emph{before} the elimination, and no regulator follows it: an
interleaved regulator placed after an elimination loses the
shaping-for-free property \cite{thomas2022worst}. The member streams
therefore keep the source bursts on their own routes, and only the last
port sees the burst of \cref{eq:frer-burst}.

\begin{algorithm2e}[t]
\caption{\textsf{Network}$(\{h\})$: bounds of all flows}
\label{alg:network}
\Input{flows with route(s), class, $b_h,T_h,J_h$}
\Output{$d_h$ for every flow, or \emph{unbounded}}
compute $\rho_x$, $l_x$, slopes (\cref{eq:slopes}) and services at every port\;
$d_{e,x}\gets0$, $\textit{ok}_{e,x}\gets$\,true for all ports\;
\Repeat{no $d_{e,x}$ or $\textit{ok}_{e,x}$ changed, or $1000$ iterations}{
	\ForEach{port $(e,x)$}{$\Sigma_{e,x}\gets0$; $\textit{fed}_{e,x}\gets$ true\;}
	\ForEach{flow $h$, each member route $e_1..e_n$}{
		$\sigma\gets b_h+\rho_hJ_h$; $\textit{b}\gets$ true\;
		\For{$i\gets1$ \KwTo $n$ (to $n-1$ if replicated)}{
			\If(\tcp*[h]{a shared $e_1$ counts once}){$i>1$ \textbf{or} first member route}{
				$\Sigma_{e_i,x}\mathrel{+}=\sigma$;
				\lIf{\textbf{not} $\textit{b}$}{$\textit{fed}_{e_i,x}\gets$ false}
			}
			\lIf{\textbf{not} $\textit{ok}_{e_i,x}$}{$\textit{b}\gets$ false}
			\lElseIf{CBS or TAS}{$\sigma\mathrel{+}=\rho_h d_{e_i,x}$}
		}
	}
	\ForEach{replicated flow $h$}{$\Sigma_{e_n,x}\mathrel{+}=$ \cref{eq:frer-burst}\;}
	\ForEach{port $(e,x)$}{
		$\textit{ok}_{e,x}\gets \textit{fed}_{e,x}\wedge\sum\rho<R_{e,x}\wedge$ load within limit\;
		$d_{e,x}\gets L_{e,x}+\Sigma_{e,x}/R_{e,x}$ if $\textit{ok}_{e,x}$ and $\le10^9$, else unbounded\;
	}
}
\lIf{not converged}{all ports unbounded}
\ForEach{flow $h$}{$d_h\gets$ \cref{eq:flow-delay} (the larger route if replicated) if all its ports are bounded, else unbounded\;}
\end{algorithm2e}

\paragraph{Fixed point.} \Cref{eq:port-delay,eq:burst,eq:frer-burst} define
the delays recursively, and \cref{alg:network} computes them by iteration
from $d\equiv0$. Larger port delays give larger bursts, hence larger
port delays, so the map is monotone and the iterates never decrease. On the star and the tree the routes are feed-forward, and the
iteration settles after at most as many passes as there are ports on the
longest route. On the ring, the clockwise routes of different flows can
form cycles of dependencies. The iteration then computes the least fixed
point of the burst equations, as in the fixed-point method for cyclic
networks \cite{leboudec2001network}. This fixed point is a valid bound: in
a network that starts empty and whose ports each serve more than the sum
of the rates of their flows, every finite nonnegative fixed point of the
burst equations bounds the bursts \cite[Thm.~2]{thomas2019cyclic}. The
first condition holds at start-up, and the second is the check of the
sustained rate (\cref{eq:port-delay}); a port that fails it is
unbounded. A port whose bound exceeds $10^9$\,\si{\micro s} is declared
unbounded, and if the iteration does not settle within 1000 passes, all
ports are. A port whose flows arrive through an unbounded port is
unbounded too (the \textit{fed} flag). With ATS, bursts do not depend on
the delays (except after an elimination, at a last port that feeds no other
port), so the cycles disappear and two or three passes suffice. The bounds
are compared with a frame-level simulation in \cref{sec:validation}.

\subsection{End-to-end bounds of a task-chain}
\label{sec:e2e}

The mapped task-chain is a sequence of segments $S_1,\dots,S_q$ separated
by communications $h_1,\dots,h_{q-1}$, inside an ECU or between ECUs
(\cref{sec:metrics}). Inside a segment, a \emph{run} is a maximal sequence
of consecutive tasks with the same period. Let $\mathrm{runs}(\xi)$ be the
runs of all segments, $\mathrm{reg}(\xi)$ the register boundaries inside
segments (consecutive tasks of different periods on the same core),
$\mathrm{cons}(\cdot)$ the task right after a boundary or a communication,
and $\mathrm{src}(h)$ the source core of $h$. The holistic sum of
\cref{sec:metrics} is
\begin{equation}
\Lambda_\xi=\sum_{r\in\mathrm{runs}(\xi)}R_{\mathrm{last}(r)}
	+\sum_{h\ \text{inter-ECU}}d_h,\label{eq:latency}
\end{equation}
and the copy-out delay of a communication $h$, counted from the end of
the WCRT of its producer, is
\begin{equation}
w_h=\begin{cases}
T^{\LET}_{\mathrm{src}(h)}+C^{\LET}_{\mathrm{src}(h)} & \text{implicit,}\\
T_h-R_{\mathrm{prod}(h)}+C^{\LET}_{\mathrm{src}(h)} & \text{strict LET.}
\end{cases}\label{eq:copy-out}
\end{equation}

\begin{proposition}[End-to-end bounds]
\label{prop:e2e}
If every task of $\xi$ meets its deadline and every inter-ECU flow has a
bound $d_h$, then
\begin{align}
\Gamma_\xi&\le\Gamma^{\mathrm{b}}_\xi=\Lambda_\xi
	+\sum_{b\in\mathrm{reg}(\xi)}T_{\mathrm{cons}(b)}
	+\sum_{h}\bigl(w_h+T_{\mathrm{cons}(h)}\bigr),\label{eq:delay}\\
\Delta_\xi&\le\Delta^{\mathrm{b}}_\xi=T_{\tau_1}+\Gamma^{\mathrm{b}}_\xi,
	\label{eq:reaction}
\end{align}
where the last sum of \cref{eq:delay} runs over all communications,
inside or between ECUs.
\end{proposition}
\begin{proof}
Follow a sample from the release of the job of $\tau_1$ that takes it.
\emph{Runs.} The tasks of a run are released together and, by (C6), run
in chain order. The data of a run is therefore complete when the job of its
last task ends, at most $R_{\mathrm{last}}$ after the common release.
\emph{Registers.} The first consumer job that reads a value written at
time $t$ is released at the first consumer release after $t$, at most
$T_{\mathrm{cons}}$ later. \emph{Communications.} With implicit
communication, an output completed at $t$ is copied at the next LET instant
of the source core, at most $T^{\LET}_{\mathrm{src}}$ later. With strict
LET, the output of a producer job released at $r$ is published at
$r+T_h$, a LET instant of the source core since $T^{\LET}_{\mathrm{src}}$
divides $T_h$; the job has ended by then, since $R_{\mathrm{prod}}\le T_h$,
and $r+T_h$ is $T_h-R_{\mathrm{prod}}$ after the bound $r+R_{\mathrm{prod}}$
counted by the run. The LET task has the highest priority, so the copy
ends at most $C^{\LET}_{\mathrm{src}}$ after the instant: $w_h$ bounds the
copy-out delay in both cases. Over the network
the message then takes at most $d_h$; its release jitter is at most
$J_h$, as assumed in its arrival curve. The consumer reads the message at
its first release after the arrival, at most $T_{\mathrm{cons}}$ later.
This release is a LET instant of the consumer core, since $T^{\LET}$
divides every period of the core, so the LET task copies the message in
at that instant, before the consumer runs. The copy-in is part of the
response time of the consumer, because the LET task interferes with it.
Summing the stages along the chain gives \cref{eq:delay}. An external
event waits at most $T_{\tau_1}$ for the next sample, which gives
\cref{eq:reaction}.
\end{proof}
The bounds use only quantities of the mapping. $\Delta^{\mathrm b}_\xi$
is the quantity compared with the deadline $D_\xi$. It bounds the maximum data age
too, since the two metrics are equal \cite{gunzel2023equivalence}.
\Cref{fig:path} shows a task-chain with two inter-ECU communications and
one register.

\begin{figure}[t]
\centering
\begin{tikzpicture}[node distance=1.4mm]
\node[task, fill=czcu] (a) {$\tau_1$};
\node[task, right=of a] (b) {$\tau_2$};
\node[draw=cmuted, dashed, rounded corners=2pt, fit=(a)(b),
	label={[note]above:ZCU2, core 9}] (s1) {};
\node[task, right=7mm of b] (c) {$\tau_3$};
\node[task, right=of c] (d) {$\tau_4$};
\node[task, right=4mm of d] (e) {$\tau_5$};
\node[draw=cmuted, dashed, rounded corners=2pt, fit=(c)(d)(e),
	label={[note]above:central computer, core 0}] (s2) {};
\node[task, right=7mm of e, fill=czcu] (f) {$\tau_6$};
\node[draw=cmuted, dashed, rounded corners=2pt, fit=(f),
	label={[note]above:ZCU4}] (s3) {};
\draw[arrow, color=corange] (s1) -- node[note, below] {$h_1$} (s2);
\draw[arrow, color=corange] (s2) -- node[note, below] {$h_2$} (s3);
\draw[arrow, dashed] (d) -- (e);
\node[note, below=4.5mm of s2, text width=8.2cm, align=left] {
runs: $\tau_1\tau_2$, $\tau_3\tau_4$, $\tau_5$, $\tau_6$;\quad
$\Lambda=R_2+R_4+R_5+R_6+d_{h_1}+d_{h_2}$\\
$\Delta^{\mathrm b}=T_1+\Lambda+T_5+(w_{h_1}+T_3)+(w_{h_2}+T_6)$};
\end{tikzpicture}
\caption{A task-chain mapped as three segments: two inter-ECU
communications, and a register between two patterns on the central
computer (dashed arrow). $\Lambda$ is the holistic sum of
\cref{eq:latency}, $\Delta^{\mathrm b}$ the bound of \cref{eq:reaction}
and $w_h$ the copy-out delay of \cref{eq:copy-out}.}
\label{fig:path}
\end{figure}

\subsection{Sensitivity: the critical scaling factor}
\label{sec:sensitivity}

The margin of a core is its critical scaling factor
\cite{lehoczky1989rate,bini2008sensitivity}: the largest $\lambda$ such that
the core stays schedulable, with its priorities, when every application
WCET $C$ becomes $\lambda C$ (the LET task unchanged). It is found by
bisection. $\lambda$ starts at 1 and doubles while the core is schedulable;
the bracket is then halved until it is within 0.1\,\%. On a core that is
not schedulable, the bisection runs between 0 and 1, so $\lambda<1$. To make small
factors meaningful, the analysis runs in nanoseconds, so that a
\SI{3}{\micro s} WCET scaled by 1.05 becomes 3150\,ns rather than staying at
\SI{3}{\micro s}. The factor of a mapping, $\lambda(\pi)$, is the smallest
over the cores it uses, since the first core to fail limits the system; it
is below 1 exactly when a core misses a deadline.
A core loaded to $\varphi$ typically has $\lambda\approx1/\varphi$.

\subsection{A lower bound on the number of cores}
\label{sec:lower-bound}

For a sub-chain $f$ and an ECU $k\in A_f$, let $u_f(k)=\min_{\theta\in
\Types(k)\cap\Theta_f}\sum_{\tau\in f}U^{\mathrm{iso}}_\tau(\theta)$ be its
utilization on the cheapest core of $k$ it may use, in isolation
(\cref{eq:wcet-iso}). Let $\Frags_k$ be the sub-chains pinned to $k$ and
$\Frags_{\mathrm{free}}$ the others.

\begin{proposition}[Lower bound]
\label{prop:lb}
For the sub-chains $\Frags$, every admissible mapping uses at least
\begin{multline}
M_{\mathrm{LB}}=\max\Bigl(\sum_{k\in\Ctrl}\ceil{\frac{U_k}{\varphi}},\\
\ceil{\frac{1}{\varphi}\Bigl(\sum_kU_k+\sum_{f\in\Frags_{\mathrm{free}}}
\min_{k\in A_f}u_f(k)\Bigr)}\Bigr)
\label{eq:lb}
\end{multline}
cores, where $U_k=\sum_{f\in\Frags_k}u_f(k)$.
\end{proposition}
\begin{proof}
Every core offers at most $\varphi$ of its own capacity (C3). A sub-chain
placed on a core of $k$ takes at least $u_f(k)$ of it, since the LET load
and the contention are non-negative. The pinned sub-chains of $k$ alone
need at least $\lceil U_k/\varphi\rceil$ cores of $k$, which gives the first
term. All sub-chains together need a total capacity of at least the
numerator of the second term, which gives the second. Constraints
(C4)--(C7) only remove mappings, so the bound holds for every admissible
mapping.
\end{proof}
The bound ignores the LET load, the schedulability gaps of non-harmonic
periods and the integrality of each zone, so it is not tight. The exact
formulation of \cref{sec:exact} gives the true optimum on smaller
instances.
